\section{Introduction}\label{sec:introduction}

Any multi-round stabilizer experiment on current quantum hardware requires a classical decoder to run between syndrome-measurement rounds, and for superconducting qubits with cycle times of roughly one microsecond, a decoder that cannot keep pace becomes the limiting factor in how long a stabilizer experiment can usefully run~\cite{fowler2012surface,terhal2015qec}. Software decoders on CPUs or GPUs typically manage a few hundred thousand decisions per second at best, three to four orders of magnitude short of that budget. Field-programmable gate arrays are the obvious answer: a shallow combinational decoder synthesised into fabric produces its output in a small, fixed number of clock cycles, and prior FPGA and ASIC decoder work has demonstrated sub-microsecond and even sub-nanosecond decode latencies at scale~\cite{das2022afs,liyanage2023scalable}.

This paper started from an attempt to build and measure exactly that: three small stabilizer-code decoder kernels in Vitis HLS, on the Xilinx Alveo U55C, with the intention of reporting hardware-measured latency and throughput. It did not go as planned, and the reason it did not go as planned is itself the paper's main finding. The kernel that was actually resynthesised for this revision reproduces its own previously reported resource footprint and critical-path timing exactly, but its previously reported pipeline latency does not reproduce under the same tool, part, and clock target (Section~\ref{sec:results}). The specific interface fix required to read a result back from the card without elevated host privileges is not free: it costs 35 to 70\(\times\) the decoder's own pipeline depth. Having identified that cost, we built the fixed kernel, placed and routed it, and verified it on a live card: it works, its self-test passes with a genuine hardware read-back, and its post-route resource and timing reports resolve a discrepancy in an earlier version of this table that previously had no real report to explain it. The decoder is close to free, correct, and now confirmed as such. The interface -- and the rest of a genuine real-time control loop -- is where the real cost and the remaining work sit.

\subsection{Scope and what this paper does not claim}\label{subsec:scope}

This paper does not claim to have built or measured a real-time QEC control loop, hardware-measured throughput, or a closed-loop demonstration with a quantum processor. It does claim a hardware-measured, hardware-verified, exhaustive correctness result for \textbf{all three} kernels (Section~\ref{subsec:res-hardware}), and a single, non-rigorous, opportunistic latency measurement for one of them (Shor) -- both explicitly scoped and caveated where they appear. Every latency and resource number in Table~\ref{tab:resources} (post-synthesis) is an HLS synthesis estimate; Table~\ref{tab:postroute} (post-route) is real for all three kernels. Every logical-error-rate curve is produced by a software model of the decoder logic, not the decoder logic executing on the FPGA fabric -- the Monte Carlo campaign and the hardware verification are separate, complementary results, not the same claim twice. Where a claim of this kind would previously have been made without a traceable artifact behind it, it has been deleted or explicitly downgraded here; a complete accounting of every claim, its evidence, and its verdict is maintained alongside this manuscript's source and referenced throughout.

What this paper does contain is a real synthesis-level characterisation of three distance-3 decoder kernels, a real, exhaustive hardware verification of all three, a real (software-only) correctness verification of all three kernels' logic against an independently written reference implementation, real Monte Carlo statistics with confidence intervals at \(10^7\) shots per point, and a quantified account of what the one architectural change needed for a working hardware read-back path actually costs -- now demonstrated, on hardware, to actually work. We believe this is a more useful contribution than an unverified real-time claim: it tells a control-stack engineer exactly where the cost of a distance-3 FPGA decoder sits, which is the question that matters when deciding whether and how to deploy one.

\subsection{Contributions}\label{subsec:contributions}

\begin{enumerate}
\item Hardware-measured, hardware-verified, \textbf{exhaustive} correctness results for all three \texttt{m\_axi}-fixed decoder kernels on a live Alveo U55C, with results read back through a genuine AXI buffer object: all $3^9=19{,}683$ Shor patterns, all $3^7=2{,}187$ Steane patterns across all three decoder modes ($6{,}561$ combinations), and the complete 16-combination Rep-3 input space, every one cross-checked against an independent software mirror with zero mismatches -- plus the first real post-route resource and timing reports in this line of work, for all three kernels.
\item A measured decomposition of an HLS decoder kernel's resource and timing cost, separating the combinational decode logic from the AXI interface, obtained by resynthesising both an AXI-Lite-only variant and an AXI-master-output variant of the same kernel and comparing them directly (Section~\ref{sec:results}).
\item A specific, quantified account of what fixing the host read-back problem costs: 35 to 70\(\times\) pipeline-depth increase for a one-element output buffer argument at the HLS-estimate stage, now confirmed on real hardware to actually deliver a working read-back path.
\item Bit-exact software mirrors of all three decoder kernels, built directly from the synthesised source (not hand-transcribed), cross-validated against an independently written reference decoder over the complete input space for the Shor kernel (262,144 of 262,144 combinations identical).
\item Monte Carlo logical-error-rate data at \(10^7\) shots per point with Wilson score intervals, including a phenomenological sweep for all three codes and a circuit-level, multi-round sweep (Stim-generated noise, PyMatching decoding) for the repetition code, which surfaces a substantial, statistically significant discrepancy with a previously reported Shor-code error rate.
\item A corrected scalability analysis of LUT-based decoding (Table~\ref{tab:lut-scaling}), with the syndrome-width convention \(m=n-k\) applied consistently, including an explicit CSS half-syndrome decomposition for the Steane code.
\item An honest account, maintained as a public, itemised ledger, of which of this paper's numbers are measured, which are estimated, which are projected, and which could not be substantiated at all.
\end{enumerate}
